%%%%%%%%%%%%%%%%%%%%%%%%%%%%%%%%%
%
%       EXERCÍCIO
%
%%%%%%%%%%%%%%%%%%%%%%%%%%%%%%%%%

\ifdefstring{\atividade}{prova}{%
    \renewcommand{\valorquestao}{\ValorQ\ pontos}
}%

\begin{exercicioBanco}[\valorquestao]
Uma variável aleatória discreta \(X\) tem a seguinte função de distribuição
%
\[
    F(x) = \left\{
    \begin{array}{rl}
        0, & \text{se } x < 0, \\
        0,1, & \text{se } 0 \le x < 1, \\
        0,4, & \text{se } 1 \le x < 3, \\
        0,6, & \text{se } 3 \le x < 4, \\
        0,85, & \text{se } 4 \le x < 10, \\
        1 & \text{se } x \ge 10. \\
    \end{array}
    \right.
\]
%
Determine:
\begin{enumerate}[(a)]
    \begin{multicols}{2}
        \item A função de probabilidade de \(X\).
        %
        \item \(P(2 \le X \le 8)\).
    \end{multicols}
\end{enumerate}
%
\end{exercicioBanco}

